\section{Conclusion}
\label{sec:conclusion}

We introduced \eps, a per-token normalized Shannon entropy over
top-$k$ log-probabilities, AST-filtered and max-aggregated into a
per-function score that also preserves per-token attribution. As a
recall-maximizing filter, the function-level claim on its own is
not distinguishing. The sequence-level mean token probability \papg
of Spiess et al.~\cite{spiess2025calibration} reads the same
logprobs, costs the same $1\times$ in API budget, and is
competitive with or better than \eps on function-level detection
for three of four OpenAI models.

The distinguishing contribution is what the function-level score
gives up to become a single scalar, and what \eps preserves instead.
A reviewer given \eps's output is not asked to re-read the function;
they are asked to evaluate the $\sim 10\%$ of semantic tokens that
cross the flag threshold. No sequence-level or inter-generational
signal produces this localization. Prior token-level work
(ASTrust~\cite{palacio2024astrust}, Vasconcelos et
al.~\cite{vasconcelos2024genprob}) shares the goal of token-level
surfacing but is presented as standalone visualization or interface
design; \eps's contribution is to package the same kind of
token-level signal as the routing layer of an end-to-end filter-plus-%
reviewer system. \papg cannot: it is one scalar per
function. Multi-sample cannot: a $0.52/0.47$ internal split
produces a unanimous output, as GPT-4o on Stripe Scenario A
empirically demonstrates ($\eps = 0.878$, sample diversity $=0.00$
across $N=5$ at $T=0.7$). UQLM detects $5\%$ of our \textsc{high}
prompts on the same benchmark. The intra-generational signal is a
different quantity from the inter-generational one, and it is the
quantity that happens to be actionable for code.

The two-stage architecture is how this becomes a working system.
Stage $1$ is the recall-maximizing filter; stage $2$ is a parallel
LLM reviewer directed to the specific high-\eps tokens rather than
asked to re-read the function. On $201$ entries across three
models, a context-primed reviewer clears $68.7\%$ of flagged entries
with zero late false positives.

The same mechanism --- early-token commitment forcing later tokens,
leaving the genuine uncertainty concentrated in a $\sim 10\%$ slice
of the generation --- plausibly applies in other structured-output
domains (SQL, configuration files, typed program synthesis);
empirical verification is left to future work. Generation is fast;
review is the bottleneck. A recall-maximizing filter, plus a
reviewer that is told exactly where to look, moves review off the
developer's critical path.
